\documentclass[../../../../SoftwareRequirementsSpecification.tex]{subfiles}
\begin{document}
% Richiede le macro definite in src/macros/useCaseMacros.tex
% (incluse nel preambolo di SoftwareRequirementsSpecification.tex)

\ucstart
\begin{xltabular}{\textwidth}{|l|X|}
  \hline
  \ucid{U-02} \\ \hline
  \ucname{Recupero Password} \\ \hline
  \ucactor{Utente (PT o Atleta)} \\ \hline
  \ucdescription{L'utente che non ricorda la propria password ne
    imposta una nuova tramite un link ricevuto via email.} \\ \hline
  \ucprecond{L'utente deve essere registrato nel sistema con
    un'email valida (vedi U-01).} \\ \hline
  \ucpostcond{L'utente ha una nuova password e può accedere con
    questa al caso d'uso U-01.} \\ \hline
  \ucmainflow{
    \item Dalla pagina di login (U-01), l'utente preme il link
      "Password dimenticata?".
    \item Il sistema mostra il form di inserimento dell'email.
    \item L'utente inserisce la propria email e preme il pulsante
      "Invia".
    \item Il sistema genera un link di reimpostazione con validità
      limitata nel tempo e lo invia all'email inserita.
    \item Il sistema mostra un messaggio che conferma l'invio
      dell'email (vedi le Note).
    \item L'utente apre l'email e preme il link ricevuto.
    \item Il sistema mostra il form di inserimento della nuova
      password.
    \item L'utente inserisce e conferma la nuova password e preme il
      pulsante "Reimposta Password".
    \item Il sistema valida la nuova password e ne sostituisce
      l'impronta a quella precedente. Il link ricevuto non è più
      valido.
    \item Il sistema mostra un messaggio di conferma e reindirizza
      l'utente alla pagina di login (U-01).
  } \\ \hline
  \ucaltflow{
    \textbf{4a. Email non registrata} \newline
    - Il sistema non invia alcuna email, ma mostra comunque il
    messaggio del passo 5 (vedi le Note). \newline
    - Il flusso termina. \newline
    \textbf{6a. Link scaduto o non valido} \newline
    - Il sistema mostra un messaggio di errore ("Link non valido o
    scaduto") con un rimando al passo 1. \newline
    - Il flusso termina. \newline
    \textbf{9a. Password non valida} \newline
    - Il sistema mostra un messaggio di errore ("Controllare i dati
    inseriti") e indica i requisiti minimi di lunghezza e
    composizione non rispettati. \newline
    - Il flusso riprende dal passo 7.
  } \\ \hline
  \ucnote{
    Il messaggio del passo 5 è mostrato indipendentemente dal fatto
    che l'email sia registrata o meno, per lo stesso motivo del
    flusso 5a di U-01: distinguere i due casi permetterebbe di
    scoprire quali indirizzi email sono registrati nel sistema.
    \newline
    Come per U-01, il sistema non memorizza la password in chiaro,
    ma solo la sua impronta calcolata con una funzione di hash.
  } \\ \hline
\end{xltabular}
\end{document}
